\section{Background and Related Work}

CODE-NN \cite{iyer2016} introduced the first fully data-driven, end-to-end neural approach to code summarization, using LSTM networks with an attention mechanism to generate natural language descriptions of C\# code snippets and SQL queries. The model was trained on a corpus automatically mined from StackOverflow, sidestepping the need for hand-labeled data. Building on this attentional encoder-decoder paradigm, code2vec \cite{alon2018} took a different angle, representing code as a fixed-length vector by decomposing it into paths through its abstract syntax tree (AST) and learning to aggregate those paths; trained on roughly 14 million methods, it was shown to predict method names accurately even on previously unseen code and to capture meaningful semantic relationships between code snippets. Code2seq \cite{Alon2018code2seqGS} extended this AST-path idea into a full sequence-generation setting, using attention over the same compositional AST paths to decode natural language sequences rather than single labels, and it outperformed prior token-sequence-based neural machine translation models across multiple languages and tasks, including code summarization and code captioning.

DeepCom \cite{hu2018} approached the problem as a structure-aware machine translation task, applying NLP techniques together with structural information extracted from Java methods to automatically generate descriptive comments, and it demonstrated clear gains in BLEU-4 scores over earlier baselines like CODE-NN by better capturing the semantics of source code rather than just its token sequence. Finally, the Transformer-based approach of Ahmad et al. \cite{ahmad2020transformerbasedapproachsourcecode} replaced recurrent architectures with a self-attention Transformer model for encoding code, finding that relative position encoding of code tokens (rather than the absolute encoding used in standard Transformers) substantially improved performance; combined with copy attention, this simple adaptation outperformed prior state-of-the-art techniques by a significant margin and has since become a common baseline architecture in the field. Together, these five papers trace the field's progression from RNN/LSTM-based sequence models, through AST-path-based structural representations, to Transformer architectures explicitly adapted to the syntactic properties of code.

Code summarization models are commonly evaluated using a combination of lexical overlap, syntactic awareness, and semantic similarity metrics. BLEU \cite{papineni2002bleu} measures the n-gram overlap between a generated summary and one or more reference summaries, emphasizing precision by rewarding exact word sequence matches. ROUGE \cite{lin2004rouge} focuses on recall and evaluates how much of the reference summary is captured by the generated text, with variants such as ROUGE-N and ROUGE-L measuring n-gram and longest common subsequence overlap, respectively. To better assess code-related text generation, CodeBLEU \cite{ren2020codebleu} extends BLEU by incorporating programming-language-specific features, including weighted n-gram matching, abstract syntax tree (AST) similarity, and data-flow consistency, providing a more comprehensive evaluation of code understanding and generation quality. In contrast, BERTScore \cite{zhang2019bertscore} leverages contextual embeddings from pre-trained transformer models to measure semantic similarity between generated and reference summaries, allowing it to capture meaning even when different words or phrasing are used. Together, these metrics provide complementary perspectives on code summarization quality, ranging from surface-level lexical agreement to deeper semantic and structural correctness.


% P3. Description of LLM-based code summaries including COBOL (Academia and industry).

% P4. Limitations of Existing approaches.

% P5. Our problem. COBOL codebase of millions of lines of code, code chunks, analytical information, and approach validation where LLMs prompts is augmented with analysis information.  Lack of ground truth. 

% P6. Approach overview. LLM-as-judge, ensemble, Token level confidence, perplexity and logit probability analysis.
